\subsection{P/PD controller}
This controller is built out of 2 blocks: predictor and traditional controller.

\subsubsection{Controllers}
The controller takes the error and performs the operations on it depending, what controller this is.
Here P and PD controllers are going to be defined. The predictors later defined are burnout, point 
mass and 2D.

\paragraph{P controller} scales the error proportionally giving
\begin{equation}
    f = k_P\cdot e
\end{equation} 

\paragraph{PD controller} scales the error proportionally and introduces damping giving
\begin{equation}
    f = k_P\cdot e + k_D \frac{de}{dt}
\end{equation} 


\subsubsection{Predictor}
The predictor is an algorithm of prediction of the final altitude $h$. The algorithm computes the apogee
for the flight given the current state of the rocket and assumption of constant air-brake extension and 
outputs the error $e$.

\paragraph{Burnout} assumes vertical only post burn dynamics. The derivation starts from equations of motion
\begin{equation}
    m \frac{dv}{dt} = -mg - kv^2.
\end{equation}

By applying approximation
\begin{equation}
    \frac{dv}{dt} = \frac{dv}{dt}\frac{dh}{dh} = \frac{dh}{dt}\frac{dv}{dh} = v \frac{dv}{dh}
\end{equation}
and $u = v^2$. One gets
\begin{equation}
    \frac{du}{dh} = -2g -\frac{2k}{m}u.
\end{equation}
This in turn can be integrated
\begin{equation}
    \int_{u}^{u_f} \frac{du}{-2g -\frac{2k}{m}u} = \int_{h}^{h_f}dh
\end{equation}
where $u = v^2$ at $h$ and $u_f = 0$ and $h_f$, gives the equation of error
\begin{equation}
    e = h_f - h = \frac{m}{2k}\ln\left(1 + \frac{k v^2}{mg}\right),
\end{equation}
where $k$ is a constant, given as
\begin{equation}
    k = \frac{\rho C_DA}{2}
\end{equation}
with all variables constant, taken as a snapshot from the current rocket state. 

\paragraph{Point mass} dynamics solver is used that is already defined within the simulation.
The simulation takes current state of the rocket and the defined environment and drag models.
The advantage of this method is that non-linear behavior can be simulated and accounts for
ideal rocket attitude change, unlike the method above.

\paragraph{2D} simulation is the improvement of point mass model, adding the real attitude dynamics.
The model introduces angle of lift and angle of attack dependece, which are going to be affected by
the brake extension.

\subsubsection{Models}
By pairing each controller with predictor model the model performances can be assessed. The table \ref{tab:1}
shows each pair and respective model ID.

\begin{table}[ht]
    \label{tab:1}
    \caption{Models made by pairing each controller with a predictor model.}
\end{table}

In theory one can also see the effect of the performance by changing the environment and drag Models
\todo{Is this needed?}
